%! Author = ben
%! Date = 27.11.2023

% Preamble
\documentclass[../main.tex]{subfiles}

% Packages

% Document
\begin{document}
	\chapter{Fachliche Kurzfassung}
	
	Für die Umsetzung der App haben wir das App-Framework Ionic verwendet.
	Dieses ermöglicht es, eine sogenannte Hybrid-App zu entwickeln.
	Eine solche Hybrid-App besteht aus einem eingebetteten Browser, welche die Benutzeroberfläche darstellt und Bindungen zu nativen Programmierschnittstellen.
	Die App verwendet für die Benutzeroberfläche die Bibliothek chakra-ui und teile des Ionic Frameworks.
	Bestimmte Teile, wie der Kalender, sind von Hand mit SCSS entwickelt.
	Im Generellen ist die verwendete Programmiersprache TypeScript, welches ein Superset der Programmiersprache JavaScript ist.
	Die App kommuniziert mit dem Backend über zwei Wege. Zum einen über eine REST-API, aber auch, mittels der Bibliothek Socket.IO, über WebSockets.
	Das Backend ist in einer Entwicklungs- und einer Produktivumgebung verfügbar, wobei der Unterschied die Umgebungsvariablen sind.
	Zum Speichern von Daten auf Seite des Backends wird MongoDB, bzw. MinIO für Bilder und andere Medien, verwendet.
	Der Backendserver ist auf Basis des Express.js Frameworks entwickelt.
	Für die Darstellung von Karten wird in der App auf Apple Maps zurückgegriffen und für die automatische Vervollständigung von Orten die Nominatim-Schnittstelle von OpenStreetMap.
	Für Inhalte, wie Artikel und Neuigkeiten wird das headless CMS directus verwendet.
	Für die Push-Benachrichtigungen wird der Dienst \dq OneSignal\dq\ verwendet.
\end{document}